%! AP Statistics — Unit 2, Lesson 2.1 — Homework KEY
\documentclass[10pt]{article}
\usepackage{apstats-article}
\usepackage{apstats-key}

%=============================================================================
\begin{document}

\pageheader{Unit 2, Lesson 2.1}{Homework --- Answer Key}
\begin{tabularx}{\linewidth}{X X X}
Name:\;\ans{Key} & Date:\; & Period:\;
\end{tabularx}

\vspace{0.1in}

\begin{tcolorbox}[colback=sky, colframe=navy, boxrule=0.8pt, arc=2mm,
    left=3mm, right=3mm, top=1.5mm, bottom=1.5mm]
\small For each news blurb: (a) identify the two variables and their types,
(b) evaluate whether the apparent association is likely \textbf{real} or possibly
\textbf{random}, and (c) identify one \textbf{confounding variable}.
\end{tcolorbox}

\vspace{0.14in}

\begin{blurbbox}[Study 1]{navylight}
\begin{headlinebox}{sky}
\small\textit{``A national study of 4{,}000 adults found that people who reported
exercising at least 3 times per week had significantly lower rates of depression
than those who exercised less.''}
\end{headlinebox}

\vspace{8pt}
\renewcommand{\arraystretch}{1.8}
\begin{tabularx}{\linewidth}{@{} l X @{}}
Two variables \& types: & \ans{Exercise frequency (categorical: 3+ times/week vs.\ less); depression rate (quantitative or categorical)} \\
Real or possibly random? & \ans{Likely real --- $n = 4{,}000$ is large; the association is unlikely to be entirely random} \\
One confounding variable: & \ans{Accept: social connection, sleep, diet, socioeconomic status, pre-existing health conditions} \\
\end{tabularx}
\end{blurbbox}

\vspace{0.14in}

\begin{blurbbox}[Study 2]{greenacc}
\begin{headlinebox}{greenbg}
\small\textit{``In a sample of 15 cities, researchers noticed that cities with more
libraries also tended to have lower crime rates.''}
\end{headlinebox}

\vspace{8pt}
\renewcommand{\arraystretch}{1.8}
\begin{tabularx}{\linewidth}{@{} l X @{}}
Two variables \& types: & \ans{Number of libraries (quantitative); crime rate (quantitative)} \\
Real or possibly random? & \ans{Possibly random --- only 15 cities; pattern may not hold in a larger sample} \\
One confounding variable: & \ans{Accept: city wealth/funding (wealthier cities can afford more libraries and also invest in crime prevention); population density; education level} \\
\end{tabularx}
\end{blurbbox}

\vspace{0.14in}

\begin{blurbbox}[Study 3]{redacc}
\begin{headlinebox}{redbg}
\small\textit{``Researchers analyzed data from 10{,}000 people and found that
individuals who sleep fewer than 6 hours per night are more likely to be overweight
than those who sleep 7--9 hours.''}
\end{headlinebox}

\vspace{8pt}
\renewcommand{\arraystretch}{1.8}
\begin{tabularx}{\linewidth}{@{} l X @{}}
Two variables \& types: & \ans{Hours of sleep per night (quantitative or categorical by category); weight status (categorical: overweight or not)} \\
Real or possibly random? & \ans{Likely real --- $n = 10{,}000$ is very large} \\
One confounding variable: & \ans{Accept: work schedules (shift workers sleep less and may eat less healthfully), stress, physical activity level, dietary habits} \\
\end{tabularx}
\end{blurbbox}

\vspace{0.14in}

\begin{reflectionbox}
\small\ans{Accept any well-reasoned response. Key elements: (1) the study is observational,
not a controlled experiment; (2) a confounding variable could explain both the explanatory
and response variable; (3) to establish causation, researchers would need to randomly assign
individuals to conditions (e.g., randomly assign some people to sleep fewer hours and observe
whether they gain weight).}
\end{reflectionbox}

\vspace{0.14in}

\begin{extensionbox}
\small
\begin{enumerate}[leftmargin=1.5em, itemsep=8pt]
  \item \ans{Accept any plausible pair. Classic examples from the website: number of
  Nicolas Cage films per year and pool drownings; per capita cheese consumption and
  deaths by tangled bedsheets.}

  \item \ans{Both variables often trend over time (increasing or decreasing year over year).
  When two unrelated trends happen simultaneously, they appear correlated even though there
  is no real relationship. This is a form of confounding by time.}

  \item \ans{Accept any reasonable two-categorical-variable research question.
  Examples: ``Is there a relationship between gender and preferred music genre?''
  or ``Do students who participate in sports also tend to participate in other clubs?''}
\end{enumerate}
\end{extensionbox}

\vspace{0.12in}
\begin{teachernote}
\textbf{Study 2 is the most important for class discussion:} Only 15 cities --- students
should flag the small sample as a reason the pattern may be random. If they miss this,
ask: ``What if we sampled a different 15 cities? Could the pattern reverse?''

\textbf{Confounding variable quality:} Push students to name confounders that are clearly
related to \emph{both} variables, not just one. ``City population'' is weak unless they
explain how it relates to both libraries and crime rate.

\textbf{Extension:} The ``time trend'' explanation is a key insight --- many spurious
correlations exist because two unrelated series both happen to trend upward (or downward)
over the same period. This previews the idea of lurking variables in regression (Lessons 2.7--2.9).
\end{teachernote}

\end{document}
